\documentclass[10pt,a4paper]{amsart}
\usepackage[margin=19mm]{geometry}
\usepackage{amsmath,amssymb,mathtools}
\usepackage[scr=rsfs,cal=euler]{mathalfa}
\usepackage{xfrac}
\usepackage[unicode,hidelinks,bookmarks=false]{hyperref}
\newcommand{\ie}{i.e.\ }
\newcommand{\cf}{cf.\ }
\newcommand{\resp}{respectively\ }
\newcommand{\Subsection}{\subsection}
\providecommand{\nobreakdash}{\nobreak\hbox{-}}
\setlength{\emergencystretch}{2em}
\multlinegap0pt
\usepackage{amssymb}
\usepackage{mathtools}
\usepackage{dsfont}
\usepackage{csquotes}
\newtheorem{thm}{Theorem}
\newtheorem{lem}{Lemma}
\title{Random Chowla's Conjecture for Rademacher Multiplicative Functions}
\author{Jake Chinis, Besfort Shala}
\date{Source: arXiv:2409.05952v2 (CC BY 4.0)}
\begin{document}
\maketitle

\noindent\emph{Source and licence.} Random Chowla's Conjecture for Rademacher Multiplicative Functions. Version arXiv:2409.05952v2 (CC BY 4.0), distributed by arXiv under the Creative Commons Attribution 4.0 International licence, http://creativecommons.org/licenses/by/4.0/, as recorded in the arXiv licence metadata for this exact version and shown on the abstract page https://arxiv.org/abs/2409.05952v2. Full licence notice: \texttt{licenses/arxiv-2409.05952-CC-BY-4.0.txt}.\par\smallskip

\noindent\textit{Editorial scope.} Counted unit: the article's thm large\_values\_theorem (source0.tex lines 208--214). The article's complete original proof of it is delivered with it (source0.tex lines 869--871, one \texttt{proof} environment of 3 lines, which the article places away from the statement). No mathematics is written, repaired or re-derived: the statement and the proof are the article's own lines, in the article's own notation, and no retained line is substituted or re-flowed.  Retained uncounted as the source-local context the proof needs: lines 859--865.  Nothing was omitted from any retained span, so the package carries no omission record.  No retained line contains a citation, so the wrapper carries no bibliography and no bibliography entry.  No retained line contains a figure environment, an image-inclusion command or drawing code, so no image is delivered and the package declares no asset dependency.  Editorial formatting only, in addition: an AMS \texttt{amsart} preamble that restates the article's own declarations and macros used by the retained lines, namely the environments \texttt{thm}, \texttt{lem} on separate counters, together with the standard packages those lines need.  The retained environments print in this document as Theorem 1, Lemma 1 in order of appearance, whereas the article prints the same items with the numbering of its own full document.  The complete native source of the article as distributed in the arXiv e-print for this exact version is bundled unchanged as \texttt{sources/165013/source0.tex}.
\begin{lem}\label{large_values_lemma}
For $X$ sufficiently large, we have 
\begin{equation}
    \max_{N\in [X, X^{(\log X)^2}]} \frac{1}{\sqrt N}\left\lvert\sum_{n\leq N} f(n^2 + 1)\right\rvert \gg \sqrt{\log\log X}\label{local_large_val}
\end{equation}
with probability $1 - O((\log\log X)^{-1/50})$.
\end{lem}

\begin{thm}
    \label{large_values_theorem} Let $f$ be a Rademacher $RMF$. There exist arbitrarily large values of $N$ such that
    \begin{equation}
        \Big|\sum_{n\leq N} f(n^2 + 1)\Big| \gg \sqrt{N\log\log N},
    \end{equation}
    almost surely.
\end{thm}

\begin{proof} [Proof that Lemma \ref{large_values_lemma} implies Theorem \ref{large_values_theorem}]
    To begin, note that \eqref{local_large_val} fails with probability $\ll (\log\log X)^{-1/50}$. Summing these probabilities over a suitably sparse sequence of values of $X$ yields a convergent series. By the Borel-Cantelli lemma, the probability that \eqref{local_large_val} fails for infinitely many of the chosen values of $X$ is $0$. Hence, there almost surely exist arbitrarily large values of $N$ for which $\frac{1}{\sqrt{N}}\Big|\sum_{n\leq N} f(n^2 + 1)\Big|\gg \sqrt{\log\log N}$, noting that $\log\log N\asymp \log \log X$ for $N\in[X, X^{(\log X)^2}$]. 
\end{proof}

\end{document}
